%! Sine and Cosine Function Graphs — Unit 6, Lesson 6.5 — Lesson Plan
\documentclass[10pt]{article}
\usepackage{precalculus-article}
\usepackage{precalculus-boxes}
\usepackage{graphicx}
\graphicspath{{images/}}
\newcommand{\TallMath}[1]{$\displaystyle #1\rule[-1.4em]{0pt}{3.2em}$}
\newcommand{\CourseName}{Precalculus}
\newcommand{\MeetingLength}{55 minutes}
\newcommand{\UnitNumberName}{Unit 6: Trigonometric Functions \quad}
\newcommand{\LessonNumberName}{Lesson 6.5: Sine and Cosine Function Graphs}

\begin{document}
\begin{center}
    {\Large\bfseries \CourseName \\
    {\small\bfseries \UnitNumberName \LessonNumberName}}
\end{center}

% ----- 1. Primary Objective -----
\begin{tcolorbox}[colback=sky, colframe=navy, boxrule=0.9pt, arc=2mm,
    left=2mm, right=2mm, top=1.5mm, bottom=1.5mm]
    \textbf{Primary Objective:} Students construct the graphs of $f(\theta)=\sin\theta$ and
    $f(\theta)=\cos\theta$ by tracking the $y$- and $x$-coordinates of a point on the unit
    circle as $\theta$ increases, then identify and describe the domain, range, period, and
    other key graphical features of each function. \quad \textit{(Big Idea: Functions)}
\end{tcolorbox}

% ----- 2. Priority Ideas & Skills -----
\renewcommand{\arraystretch}{1.6}
\begin{skillbox}[Priority Ideas \& Skills]{goldbox}
\begin{minipage}[t]{0.47\linewidth}\vspace{0pt}
    \textbf{AP Skills}\\[4pt]
    \textbf{Skill 2.A} --- Identify mathematical characteristics of a function from a
    graphical representation (e.g.\ domain, range, period, zeros, max/min).\\[4pt]
    \textbf{Skill 3.A} --- Describe the characteristics of a function using appropriate
    mathematical language (oscillation, amplitude, period, midline).\\[4pt]
    \textbf{LO 3.4.A} --- Construct representations of the sine and cosine functions
    using the unit circle.
\end{minipage}
\hfill
\begin{minipage}[t]{0.47\linewidth}\vspace{0pt}
    \textbf{Key Understandings (EKs)}\\[4pt]
    \begin{itemize}[leftmargin=1.2em, itemsep=3pt]
        \item The $y$-coordinate of the unit-circle point at angle $\theta$ gives
              $f(\theta)=\sin\theta$; domain is all reals. (EK 3.4.A.1)
        \item As $\theta$ increases, $\sin\theta$ oscillates between $-1$ and $1$,
              tracing vertical displacement. (EK 3.4.A.2)
        \item The $x$-coordinate gives $f(\theta)=\cos\theta$; domain is all reals.
              (EK 3.4.A.3)
        \item As $\theta$ increases, $\cos\theta$ oscillates between $-1$ and $1$,
              tracing horizontal displacement. (EK 3.4.A.4)
    \end{itemize}
\end{minipage}
\end{skillbox}

% ----- 3. Vocabulary -----
\begin{skillbox}[{Vocabulary, Concepts, \& Theorems}]{greenbox}
\renewcommand{\arraystretch}{1.45}
\begin{tabularx}{\linewidth}{@{} l X @{}}
    \textbf{Sine function}        & $f(\theta)=\sin\theta$; the $y$-coordinate of the
                                    unit-circle point at angle $\theta$. \\
    \textbf{Cosine function}      & $f(\theta)=\cos\theta$; the $x$-coordinate of the
                                    unit-circle point at angle $\theta$. \\
    \textbf{Oscillation}          & Repeated back-and-forth variation between two extreme
                                    values. \\
    \textbf{Amplitude (informal)} & Half the total vertical range; equals $1$ for both
                                    $\sin\theta$ and $\cos\theta$. \\
    \textbf{Period}               & Smallest positive input change for one complete cycle;
                                    $2\pi$ for both functions. \\
    \textbf{Midline}              & Horizontal line halfway between max and min; $y=0$ for
                                    both functions. \\
    \textbf{Maximum}              & Greatest output value; $1$ for both functions. \\
    \textbf{Minimum}              & Least output value; $-1$ for both functions. \\
    \textbf{Zero (of a function)} & An input $\theta$ where $f(\theta)=0$. \\
\end{tabularx}
\end{skillbox}

% ----- 4. Activate Prior Knowledge -----
\begin{skillbox}[Activate Prior Knowledge \& Spiral Review]{sky}
\begin{tabularx}{\linewidth}{@{}X c@{}}
    \begin{minipage}[t]{0.58\linewidth}\vspace{0pt}
        \textbf{Students should already be able to:}
        \begin{itemize}[leftmargin=1.2em, itemsep=3pt]
            \item State exact values of $\sin\theta$ and $\cos\theta$ at the standard
                  unit-circle angles ($0,\,\tfrac{\pi}{6},\,\tfrac{\pi}{4},\,\tfrac{\pi}{3},
                  \,\tfrac{\pi}{2},\,\pi,\,\tfrac{3\pi}{2},\,2\pi$) --- from Lesson 3.3.
            \item Recognize a periodic function from a description --- from Lesson 3.1.
            \item Read and interpret function values from a Cartesian graph.
        \end{itemize}
        \vspace{4pt}
        \textit{The warm-up recalls Lesson 3.3 exact values and asks whether $\sin\theta$ is
        increasing on $[0,\tfrac{\pi}{2}]$, motivating today's graphical view.}
    \end{minipage}
    &
    \begin{minipage}[t]{0.35\linewidth}\vspace{0pt}\centering
        \small\textit{(Warm-up uses unit-circle\\exact values from L3.3.)}
    \end{minipage}
\end{tabularx}
\end{skillbox}

% ----- 5. Hook -----
\begin{skillbox}[Hook]{sky}
    \textbf{Entry Question (Think-Pair-Share, 90 sec):} ``We know the $y$-coordinate of a
    point on the unit circle at every angle. What does the graph of that $y$-coordinate
    \emph{versus the angle} look like?''

    \vspace{4pt}
    Expect answers like ``it goes up then down'' or ``maybe a wave.'' After sharing, ask the
    same for the $x$-coordinate to preview $\cos\theta$. Connect to the periodic phenomena
    from Lesson 3.1 (tides, temperature cycles): ``Today we find the exact mathematical
    function behind those repeating patterns.''
\end{skillbox}

% ----- 6. Lesson -----
\begin{skillbox}[Lesson]{sky}
\begin{multicols}{2}
    \textbf{Building the Graph of $\boldsymbol{f(\theta)=\sin\theta}$}

    \vspace{4pt}
    Have students fill in a table of $(\theta,\sin\theta)$ pairs using their unit-circle
    reference:

    \vspace{4pt}
    \renewcommand{\arraystretch}{1.3}
    \begin{tabular}{@{}c|c@{}}
        $\theta$ & $\sin\theta$ \\\hline
        $0$               & $0$ \\
        $\tfrac{\pi}{6}$  & $\tfrac{1}{2}$ \\
        $\tfrac{\pi}{4}$  & $\tfrac{\sqrt{2}}{2}$ \\
        $\tfrac{\pi}{3}$  & $\tfrac{\sqrt{3}}{2}$ \\
        $\tfrac{\pi}{2}$  & $1$ \\
        $\pi$             & $0$ \\
        $\tfrac{3\pi}{2}$ & $-1$ \\
        $2\pi$            & $0$ \\
    \end{tabular}

    \vspace{6pt}
    \textbf{Pose:} ``As $\theta$ increases from $0$ to $\tfrac{\pi}{2}$, what happens to
    $\sin\theta$? From $\tfrac{\pi}{2}$ to $\pi$?''

    Plot the ordered pairs on axes with $\theta$ horizontal ($0$ to $2\pi$) and $y$ vertical
    ($-1$ to $1$). Connect smoothly --- the \textbf{sine wave} emerges.

    \columnbreak

    \textbf{Building the Graph of $\boldsymbol{f(\theta)=\cos\theta}$}

    \vspace{4pt}
    Repeat with $x$-coordinates (horizontal displacement):

    \vspace{4pt}
    \begin{tabular}{@{}c|c@{}}
        $\theta$ & $\cos\theta$ \\\hline
        $0$               & $1$ \\
        $\tfrac{\pi}{2}$  & $0$ \\
        $\pi$             & $-1$ \\
        $\tfrac{3\pi}{2}$ & $0$ \\
        $2\pi$            & $1$ \\
    \end{tabular}

    \vspace{6pt}
    \textbf{Pose:} ``How does the cosine graph compare to the sine graph in shape? In
    starting value?''

    \vspace{6pt}
    \textbf{Key Features --- whole-class debrief}

    \vspace{4pt}
    From the graphs, elicit:
    \begin{itemize}[leftmargin=1.2em, itemsep=2pt]
        \item Domain: all real numbers
        \item Range: $[-1,\,1]$
        \item Period: $2\pi$
        \item Midline: $y=0$
        \item Maximum: $1$; Minimum: $-1$
        \item Zeros of $\sin\theta$ on $[0,2\pi)$: $0,\,\pi,\,2\pi$
        \item Zeros of $\cos\theta$ on $[0,2\pi)$: $\tfrac{\pi}{2},\,\tfrac{3\pi}{2}$
    \end{itemize}
\end{multicols}
\end{skillbox}

% ----- 7. Active Monitoring -----
\begin{skillbox}[Active Monitoring]{redbox}
    \textbf{Circulate and watch for:}
    \begin{itemize}[leftmargin=1.2em, itemsep=3pt]
        \item Confusion about which coordinate maps to sine vs.\ cosine --- redirect to the
              unit-circle diagram.
        \item Students plotting $(\sin\theta,\,\theta)$ instead of $(\theta,\,\sin\theta)$ ---
              ask: ``Which quantity is the input? Which is the output?''
        \item Students who stop at $2\pi$ and don't see the periodic extension --- ask:
              ``What happens after one full revolution?''
    \end{itemize}
    \textbf{Cold-call prompts:}
    \begin{itemize}[leftmargin=1.2em, itemsep=3pt]
        \item ``At $\theta=\tfrac{3\pi}{2}$, is $\sin\theta$ positive, negative, or zero?
              How does the unit circle show that?''
        \item ``Why do both graphs have the same range $[-1,1]$?''
        \item ``Where does $\cos\theta=1$? Name all solutions in $[0,4\pi]$.''
    \end{itemize}
\end{skillbox}

% ----- 8. Group Work & Differentiation -----
\begin{skillbox}[Group Work \& Differentiation]{redbox}
\begin{multicols}{3}
    \textbf{Tier R --- Remediate}\\
    Given a complete table of $(\theta,\sin\theta)$ values for key angles, answer: where are
    the zeros? Where is the maximum? On which interval from $0$ to $2\pi$ is $\sin\theta$
    increasing? What is the period? Use the unit-circle diagram as a reference.

    \columnbreak
    \textbf{Tier A --- Approaching}\\
    Given a partial table of $(\theta,\cos\theta)$ values with some blanks, complete the table
    using the unit circle or symmetry. Then answer: on what intervals in $[0,2\pi]$ is
    $\cos\theta$ positive? negative? What transformation would map the sine graph to the cosine
    graph?

    \columnbreak
    \textbf{Tier E --- Extension}\\
    Use table values to estimate the average rate of change of $\sin\theta$ on
    $[0,\tfrac{\pi}{4}]$, $[\tfrac{\pi}{4},\tfrac{\pi}{2}]$, and
    $[\tfrac{\pi}{2},\tfrac{3\pi}{4}]$. Describe the pattern and explain it using the
    unit circle.
\end{multicols}
\end{skillbox}

% ----- 9. Individual Work & Assessment -----
\begin{skillbox}[Individual Work \& Assessment]{redbox}
    \textbf{Exit Ticket (2 items, 5 min):}
    \begin{enumerate}[leftmargin=1.5em, itemsep=4pt]
        \item State the domain, range, and period of $f(\theta)=\cos\theta$.
        \item A student says ``$\sin\theta$ and $\cos\theta$ have the same graph.''
              Correct the student's claim with one specific example showing a difference.
    \end{enumerate}
    \vspace{4pt}
    \textbf{Assessment note:} Collect exit tickets as students leave. Sort into three piles
    (both correct / one correct / neither) to plan Tier R support in Lesson 3.5.
\end{skillbox}

% ----- 10. Reinforcement & Extension -----
\begin{skillbox}[Reinforcement \& Extension]{goldbox}
    \textbf{Homework:} 5--6 problems reinforcing domain, range, period, and zeros of
    $\sin\theta$ and $\cos\theta$; includes True/False with justification, a conceptual
    question on periodicity, and an AP-style MC. Optional extension on concavity.

    \vspace{4pt}
    \textbf{Extension:} On what intervals is $\cos\theta$ concave down? Connect the answer
    to the curvature of the unit-circle arc.

    \vspace{4pt}
    \textbf{Preview --- Lesson 3.5:} In the next lesson we transform these graphs by changing
    amplitude (vertical stretch), period (horizontal stretch), and midline (vertical shift) to
    model real-world sinusoidal phenomena.
\end{skillbox}
\end{document}
